%! Complex Numbers — Unit 3, Lesson 3.4 — Warm-Up (Answer Key)
\documentclass[10pt]{article}
\usepackage{algebra2-article}
\usepackage{algebra2-key}   % pulls in -boxes; do NOT also load -boxes
\newcommand{\TallMath}[1]{$\displaystyle #1\rule[-1.4em]{0pt}{3.2em}$}

\begin{document}

\pageheader{Unit 3, Lesson 3.4}{Warm-Up --- Answer Key}
\namedateperiod

\vspace{0.10in}

\begin{notesbox}{Getting Ready: Skills We'll Use Today}
\small These four skills power everything about complex numbers. Work quickly.

\vspace{4pt}
\textbf{1. Simplify each radical.} Pull out the largest perfect square (simplest radical form).
\[
  \sqrt{50}=\ans{$5\sqrt{2}$} \qquad \sqrt{18}=\ans{$3\sqrt{2}$} \qquad
  \sqrt{48}=\ans{$4\sqrt{3}$} \qquad \sqrt{8}=\ans{$2\sqrt{2}$}
\]

\vspace{4pt}
\textbf{2. Multiply the binomials.} Distribute (FOIL), then combine like terms.
\[
  (x+3)(x-3)=\ans{$x^2-9$} \qquad (2x+1)(x-4)=\ans{$2x^2-7x-4$}
\]

\vspace{4pt}
\textbf{3. Combine like terms.} Add or subtract; watch the signs.
\[
  (3x+5)+(2x-8)=\ans{$5x-3$} \qquad (7-4x)-(2-9x)=\ans{$5x+5$}
\]

\vspace{4pt}
\textbf{4. Hit the wall.} From Lesson 3.3: solve $x^2=-4$ for a \emph{real} number $x$.
\begin{itemize}\setlength{\itemsep}{1pt}
  \item Is there a real number whose square is $-4$? \ans{No --- any real number squared is $\ge 0$.}
  \item So how many \emph{real} solutions does $x^2=-4$ have? \ans{None (no real solution).}
\end{itemize}
\emph{Today we build a new kind of number so that this equation \textbf{can} be solved.}
\end{notesbox}

\begin{teachernote}
Every item is a direct prerequisite. Item~1 keeps radicals in simplest form --- exactly what
$\sqrt{-50}=5i\sqrt2$ will need once the negative comes out as $i$. Item~2 is FOIL, the engine for
multiplying complex numbers (and $(x+3)(x-3)=x^2-9$ previews the conjugate product). Item~3 is
combining like terms --- adding/subtracting complex numbers is nothing more than this with a real
part and an ``$i$ part.'' Item~4 reopens the Lesson~3.3 wall ($x^2=-4$) that today's imaginary unit
$i$ finally breaks. Keep item~4 on the board to launch the hook.
\end{teachernote}

\end{document}
